\section{数据集与任务定义}

本章我们依次介绍实验所用的数据集、预处理流程、数据划分策略、任务设定以及评价指标。

\subsection{数据集介绍}

本次课程设计采用 Thought2Text 数据集。该数据集从 6 名健康被试处采集得到，每名被试观看 40 个 ImageNet 类别下的自然图像，同步通过 128 通道 EEG 设备记录脑电信号。整个数据集中约有 11965 条 EEG trial，对应大约 1996 张不同的图像。经过降采样后，每条 EEG 数据的维度为 $[64, 250]$（64 通道，250 个时间点）。图像统一缩放到 $224\times224$ 作为模型输入。

需要说明的是，Thought2Text 是一个 EEG 与图像配对的数据集，原始标注为类别标签而非自然语言描述，因此其直接用途是 EEG-image 类别语义学习。针对生成式任务的需求，我们借助预训练 BLIP-base 模型为每张图像生成伪描述文本，作为后续训练生成模型的目标。

数据集概况见表~\ref{tab:dataset}。

\begin{table}[htbp]
  \centering
  \caption{Thought2Text 数据集基本信息}
  \label{tab:dataset}
  \begin{tabular}{l c}
    \toprule
    项目 & 数值 \\
    \midrule
    图像总数 & 约 1996 张 \\
    EEG trial 总数 & 约 11965 条 \\
    类别数 & 40 个 \\
    被试数 & 6 名 \\
    EEG 形状 & $64\times250$ \\
    图像输入分辨率 & $224\times224$ \\
    原始 EEG 通道数 & 128 \\
    \bottomrule
  \end{tabular}
\end{table}

\subsection{数据预处理}

数据预处理分为 EEG 信号处理、图像处理和文本处理三个部分。

\textbf{EEG 预处理}：原始 128 通道 EEG 降采样为 64 通道，逐通道做 z-score 标准化。所有 trial 统一截取为 250 个采样点。为了减轻训练时的磁盘读取压力，我们将训练集、验证集和测试集的全部 EEG 数据预先存为 numpy 数组格式，训练过程中直接从内存加载。

\textbf{图像处理}：统一使用 CLIP ViT-B/32 的标准预处理流程（缩放到 $224\times224$、CenterCrop、ImageNet 标准化）。为了节省训练时的计算量，我们将所有图像的 CLIP 视觉特征（包括 pooled embedding 和 patch tokens）预先计算并以 float16 精度存为 .npy 文件，训练时读入再转成 bfloat16 使用。

\textbf{文本处理}：为 40 个类别分别构造 ``a photo of a \{class\_name\}'' 形式的 prompt，用冻结的 CLIP text encoder 提取 512 维文本原型向量，预计算后固定存储，训练和评估阶段不再更新。

\textbf{视觉退化生成}：六种退化条件在数据加载时在线生成，不需要预先存储退化图像。退化操作基于 torchvision 的标准算子实现，在 DataLoader 的 worker 进程中执行。

\subsection{数据划分}

数据划分采用 image-level 方案，即按图像 ID 而非 trial ID 进行切分。这样设计的出发点是为了避免信息泄漏：如果按 trial 切分，同一张图像的不同 trial 可能分属训练集和测试集，模型有可能通过记忆图像本身的视觉特征来完成任务，而不是真正利用 EEG 中的语义信息。采用 image-level split 之后，训练集和测试集之间不存在共享图像。

划分的具体数量详见表~\ref{tab:split}。验证集用于模型选择和早停判断，测试集只在最终评估时使用。

\begin{table}[htbp]
  \centering
  \caption{数据集训练/验证/测试划分}
  \label{tab:split}
  \begin{tabular}{l c c c}
    \toprule
    划分 & 图像数 & EEG trial 数 & 类别数 \\
    \midrule
    训练集 & 1330 & 7970 & 40 \\
    验证集 & 333 & 1998 & 40 \\
    测试集 & 333 & 1997 & 40 \\
    \bottomrule
  \end{tabular}
\end{table}

\subsection{任务定义}

我们围绕两个核心任务展开实验。

\textbf{任务一：EEG 辅助语义预测}。输入退化图像及对应的 EEG trial，模型输出 40 个类别上的概率分布。该任务旨在检验真实配对的 EEG 信号是否能在图像退化条件下帮助模型做出更准确的类别判断。评估阶段设置了五种输入模式：纯视觉（vision-only）、真实 EEG（real EEG）、shuffled EEG（打乱脑电）、随机 EEG（random EEG）和仅用 EEG（eeg-only）。

\textbf{任务二：生成式图像描述}。输入退化图像及对应的 EEG trial，模型生成对应的自然语言描述。该任务考察 EEG 增强后的视觉表示是否能让生成式视觉语言模型输出比纯视觉更准确的描述。评估阶段设置了四种模式：vision-only、real EEG、shuffled EEG 和 random EEG。

\subsection{评价指标}

\textbf{语义预测指标}：

\begin{itemize}
  \item \textbf{Top-1 Accuracy}：模型预测的最高概率类别与真实类别一致的比例；
  \item \textbf{Top-5 Accuracy}：真实类别位于预测前 5 类别中的比例；
  \item \textbf{Class Hit}：通过 CLIP 文本原型相似度进行零样本类别判定，命中真实类别的比率；
  \item \textbf{Real-Vision Gap}：real EEG 与 vision-only 的指标差值，正值表示 EEG 提供了正向辅助；
  \item \textbf{Real-Shuffled/Random Gap}：real EEG 与 shuffled/random EEG 的差值，是判断 EEG 效果是否来自真实配对的关键对照指标；
  \item \textbf{Win Rate}：在多种退化条件或随机种子中，real EEG 超过对照组的比例。
\end{itemize}

\textbf{生成式指标}：

\begin{itemize}
  \item \textbf{有效输出率（Valid Caption Rate）}：生成文本中合法有效描述的比例。有效描述定义为非空、非纯乱码、非代码片段和非 URL 的文本；
  \item \textbf{无效输出率（Invalid Output Rate）}：生成文本无效的比例；
  \item \textbf{Caption Class-Hit}：对生成描述进行 CLIP 零样本类别判定，命中真实类别的比率；
  \item \textbf{Distinct Caption Count}：生成的不同描述文本总数，反映输出多样性；
  \item \textbf{平均描述长度（Avg Caption Length）}：生成描述的平均 token 数。
\end{itemize}
